\chapter{奶蛙进校园：空间实践、组织形态与低强度协商}\label{ch:10}

\begin{nlkeypoints}
  \begin{enumerate}[label=(\arabic*),leftmargin=2.4em]
    \item 五校案例（哈工大、上交、海大、北航、浙大，均 \conf{C} 级）
          在时间上呈\keyw{扩散序列}：从单一立牌亮相（哈工大）
          到流动实践（上交）、复数立牌（海大）、
          组织化同好会（北航）、官方物料采纳（浙大）——
          与创新扩散的采用者类别在形态上相容
          （\conf{C} 级案例 + \conf{M} 级结构归位）。
    \item 立牌是奶蛙在校园的\keyw{物化符号}：
          它把线上迷因固定为可合影、可搬运、可巡游的实物，
          使符号实践进入物理空间并接受校方管理——
          2026-09-28 北航“9 小时立牌”事件展示了
          \keyw{移除—纪念—再创作}的低强度协商回路。
    \item 本章给出\keyw{协商强度谱}（\conf{M} 级）：
          按“校方反应—学生反应—是否留下持久物”三维
          给案例分级；北航事件的三维读数表明，
          校园协商以“边界重申 + 符号存续”收尾，
          未见冲突性对抗。
  \end{enumerate}
\end{nlkeypoints}

前九章的分析都发生在“线上”——平台、语料、生产链。
本章进入物理空间：奶蛙如何出现在毕业典礼、绿园、
招生宣传与同好会群聊里。材料全部来自第一版教材的
高校活动记录（\conf{C} 级），单一来源、无法独立印证，
这是本章证据地位的底限，分析结论全部标注 \conf{M} 级。
本章回答问题 Q3 的最后一环：当迷因从信息流走进校园，
它的传播结构、组织载体与权威关系发生了什么变化。

\section{案例库：五校事件整理}\label{sec:10:cases}

\begin{table}[htbp]\centering
  \caption{高校奶蛙事件整理（据第一版教材，全部 \conf{C} 级；
          “类型”列为本书归类，\conf{M} 级）}
  \label{tab:10:cases}
  \small
  \begin{tabularx}{\textwidth}{@{}lllX@{}}
    \toprule
    学校 & 时间 & 类型 & 事件 \\ \midrule
    哈尔滨工业大学 & 2026-06 & 仪式亮相 &
      毕业典礼暨学位授予仪式出现第一个广为人知的奶蛙立牌，
      被视为各校“立牌争先赶超”阶段的开启 \\
    上海交通大学 & 2026 年 & 流动实践 &
      率先发起“流动奶蛙”：把立牌带到校园各处制造偶遇；
      最早的流动奶蛙至今仍在使用 \\
    中国海洋大学 & 2026 年 & 复数立牌 &
      “奶 OUC”四奶蛙立牌，第一个广为人知的复数立牌 \\
    北京航空航天大学 & 2026-06-24 起 & 组织化 &
      奶龙同好会成立；2026-08 初破 100 人；
      26 级新生入学后于 2026-10 初破 200 人 \\
    北京航空航天大学 & 2026-09-28 & 空间冲突 &
      学院路校区绿园投放流动奶蛙立牌，旁现丰川祥子立牌；
      当晚被保安带走，存活约 9 小时 \\
    浙江大学 & 2026 年 & 官方采纳 &
      某地招生宣传立牌中出现浙大奶蛙形象 \\
    \bottomrule
  \end{tabularx}
\end{table}

表 \ref{tab:10:cases} 的“类型”列暗含一个
时间—结构假设（\conf{M} 级）：五个案例恰好覆盖了
创新扩散理论中采用者的经典类别序列——
\term{创新者}（哈工大：首个、仪式场合、高风险亮相）、
\term{早期采用者}（上交：发明新玩法“流动奶蛙”）、
\term{早期大多数}（海大：把单点实践扩展为复数形态）、
\term{后期大多数}（各校同好会的组织化复制，
北航为其显例）、以及\keyw{守门人采纳}
（浙大：官方物料出现该形象）\cite{rogers2003diffusion}。
必须强调两点：其一，这是\keyw{形态上的相容}，
不是对扩散过程的实证检验——案例时间粒度太粗
（精确到月），且五校样本是“出圈案例”的截尾样本
（见 \ref{sec:10:network} 的现存样本偏差讨论）；
其二，序列若成立，它预测 2026 年下半年还会有
更多“守门人采纳”型事件（官方物料化），
这是一个可跟踪的预测。

\section{空间实践：立牌、流动与合影}\label{sec:10:space}

\subsection{立牌：从梗图到实物}

立牌是奶蛙符号在物理空间的锚点。第六章指出，
“摆架势”是衍生动作中“从形象动作迁移为使用者
身体姿态的通道”；立牌把这一通道反过来用——
\keyw{把形象固定为可合影的实物}，使身体程式
（摆架势、合影）有了一个空间坐标。
用 Jenkins 等人的语言说，立牌是可传播性的
\keyw{物化形态}：内容不依赖特定载体传播
（命题 \ref{prop:09:unit}），
而立牌让“载体”本身可以被搬运、被巡游、被拍下再传回线上
\cite{jenkins2013spreadable}。

\subsection{流动奶蛙：制造偶遇}

上交的“流动奶蛙”把立牌从静物变为\keyw{移动事件}：
“把奶蛙带到学校的各个角落，在任意一处地方偶遇惊喜”
（\conf{C}，第一版教材口径）。
这一实践的文化逻辑是\keyw{偶遇的仪式化}——
把概率性事件（碰巧刷到梗）改写为空间性事件
（碰巧遇见蛙）。Turner 的阈限理论提供了一种读法：
毕业典礼、校园绿地本是日常结构的通道，
立牌的闯入制造了一个\keyw{微阈限}时刻——
参与者短暂脱离日常角色（学生、行人），
进入“与一只大笑青蛙合影”的反结构状态，
合影与转发则是阈限体验的\keyw{存档}
\cite{turner1969ritual}。
这一读法是本书的诠释（\conf{M} 级），
它的可检验推论是：流动奶蛙的合影率
（遇见者拍照比例）应显著高于固定装置的合影率——
因为“偶遇”比“陈设”更接近阈限体验的触发条件。

\subsection{FoMO 与纳新}

第二章预言了“准入测试在纳新中反复出现”
（\ref{sec:02:meta}），北航同好会的增长提供了
检验材料：2026-06-24 成立，8 月初破 100 人，
26 级新生入学后 10 月初破 200 人
（\conf{C}，第一版教材口径）。
新生入学与会员数翻倍的时间重合提示
\keyw{错失恐惧}（\en{FoMO}）的动因结构：
FoMO 指对他人正在经历自己缺席的有益经验的
普遍性忧虑，它驱动社交媒介使用与从众行为
\cite{przybylski2013fomo}。
在纳新语境下，刷屏的立牌照片与同好会群聊
构成“大家都在里面”的可见证据，
FoMO 把兴趣转化为入会动作
（\conf{M} 级分析；第十一章问卷含 FoMO 量表
以便检验）。

\section{组织形态：同好会}\label{sec:10:org}

第一版教材记录的北航奶龙同好会成立目标
（\conf{C}）是：“让原飞梦社的抽象老登们
有个神圣的水群”。“水群”（社群用语）指低强度的
日常闲聊；记录同时说明该会因“奶蛙前所未有的强横热度”
而人数激增，“成为飞梦之后北航校园文化新的中坚力量之一”
\conf{C}。

这条记录的学理价值在于它显示了\keyw{组织的双轨性}：
同好会的名义功能（水群）与实际功能（校园迷因的
组织载体：纳新、贺图、表情包生产、立牌活动）
是分开的。Hebdige 笔下的亚文化很少有正式组织，
奶蛙的同好会则是\keyw{半正式}的——
有成立日期、有人数里程碑（100 人贺图、200 人贺图）、
有名称与身份，但无正式社团注册的记载
（\conf{C} 级材料未提及注册状态，本书不推断）。
第二章“准入测试”的功能分析在此落位：
贺图与校徽奶蛙表情包（“数分奶蛙”“捧着航概的奶蛙”
\conf{C}）是\keyw{校本化的迷因资产}，
它们把全国性梗转写为本校身份标记——
持有并使用校本奶蛙表情包，
即是“本校圈内人”的默示声明。

\section{协商：9 小时立牌事件与协商强度谱}\label{sec:10:negotiation}

\subsection{事件经过}

据第一版教材记录：2026 年 9 月 28 日，
一个匿名者在北航学院路校区绿园投放流动奶蛙立牌，
旁边随后出现丰川祥子立牌；学院路学生纷纷与之合影；
当晚立牌被保安发现并带走，存活约 9 小时；
同好会“对此深表遗憾”，成员自发在群内刷屏一段
化用宗教语体的赞诗（“我的肉化作黄桃罐头……
为使世界获得笑容”，\conf{C}，第一版教材记录；
本书只作语体分析，不复制全文）。

\subsection{低强度协商的三维读数}

事实手册指出该事件的关键特征：
官方（保安/校方）与自组织（同好会）之间
\keyw{并未发生冲突性对抗}，而是以
“移除—纪念—再创作”完成一轮低强度协商。
本书把这一观察升级为可分级的教学框架：

\begin{table}[htbp]\centering
  \caption{立牌活动的协商强度谱（本书构造的教学模型，\conf{M} 级；
          案例读数基于 \conf{C} 级材料）}
  \label{tab:10:spectrum}
  \small
  \begin{tabularx}{\textwidth}{@{}lXXX@{}}
    \toprule
    维度 & 低强度（0） & 中强度（1） & 高强度（2） \\ \midrule
    校方反应 & 例行移除、未追责 &
      公开通告、约谈组织者 &
      纪律处分、定性为违规活动 \\
    学生反应 & 自嘲、默默再投放 &
      纪念文本、纪念二创 &
      组织化抗议、对抗性表态 \\
    持久物 & 无 & 线上纪念物 &
      制度化安排（如申请合规展位） \\ \bottomrule
  \end{tabularx}
\end{table}

北航事件的读数：校方反应 0（例行移除，
无公开追责记录）；学生反应 1（纪念文本与刷屏，
无对抗）；持久物 1（线上纪念物；且“流动奶蛙”
实践本身在事件后延续——最早的流动奶蛙仍在使用）。
合计 2，属\keyw{低—中强度区间}。
这个分级的用途不是给事件打分，而是让\keyw{比较}成为可能：
不同学校、不同时间的立牌事件可以放在同一谱上对照，
“协商强度随什么变量变化”（学校管理风格？
投放场合的正式程度？立牌内容？）由此变成可提问的经验问题。

\subsection{赞诗的语体分析}

赞诗值得单独一读，因为它是第六章“抽象文化”
定义中\keyw{对严肃语体的戏仿}的最佳样本：
它化用宗教圣体文的句式（“我的肉……谁若吃了……
必要……直到永远”），把一只橡皮鸭子式的玩偶立牌
写进最庄重的语体框架。语体错配在这里承担三件事：
其一，\keyw{哀悼的合法化}——为一个玩偶立牌“致哀”
本会尴尬，宗教语体的庄重性允许成员
以夸张完成真诚；其二，\keyw{圈内测试}——
识别这段文本的互文来源需要特定文化库存，
读懂即认证；其三，\keyw{软性表态}——
刷屏行为本身向“带走立牌的管理方”传递了
不满，但语体的戏谑性使该表态随时可撤回
（“我们只是在玩梗”）。
这与第七章 M2（群体稀释）的“印象缓冲”结构同源：
\keyw{用可撤回的语体，说不可撤回的态度}。

\begin{nlcritique}[批判性反思：管理方的视角与道德恐慌的风险]
  本章的案例都以学生视角记录（来源即学生教材），
  校方视角完全缺席。公平起见必须列出管理方的
  合理关切：未报备的立体装置涉及公共空间秩序、
  消防与通行安全、以及“下一个投放者放什么”
  的滑坡问题（丰川祥子立牌的出现已经演示了
  跟随投放的即兴性）。同时要警惕反向的过度反应：
  若把校园玩梗上升为“不良文化”进行运动式清理，
  即落入 Cohen 所述\keyw{道德恐慌}的框架——
  把一个低危害的风格现象建构成对秩序的威胁
  \cite{cohen1972folkdevils}。
  协商强度谱的意义正在于提供第三条路：
  双方都以低强度动作表达立场（移除对纪念），
  边界被重申（不经报备的装置会被移走），
  符号实践得以延续（流动奶蛙照常）。
  \keyw{低强度协商不是冲突的失败，而是冲突的替代}。
\end{nlcritique}

\section{传播网络：校园作为高密度混合网络}\label{sec:10:network}

第七章 M4 预测“线下共同在场密度延长素材寿命”
（假设 \ref{hyp:07:m4}），并承诺校园传播网络的结构分析。
本节给出框架（分析框架 \conf{M} 级）。

\subsection{弱连接与强连接的分工}

Granovetter 的弱连接优势命题：新信息更可能
经由弱连接（泛泛之交）到达，因为强连接圈子内部
信息早已同质\cite{granovetter1973weak}。
校园网络是罕见的\keyw{高密度混合}结构：
宿舍与班级提供强连接（复读、互发表情包，
强化趣味、延迟疲劳，对应 M4），
跨院系社团、校园广场与信息流提供弱连接
（新变体、新梗的输入，对应第七章“新变体重新激活”
的 $\omega$ 项）。预测：\keyw{一个校园迷因的寿命
取决于强连接复读维持与弱连接变异供给的乘积}——
只有强连接而无新变体，趣味在复读中缓慢耗尽；
只有弱连接而无复读，梗无法沉淀为本校身份标记。
校本表情包（数分奶蛙、航概奶蛙）正是
弱连接输入被强连接复读固化后的\keyw{沉淀物}。

\subsection{多群 SIR 与现存样本偏差}

第五章批判框承认单群 SIR 模型的第一处局限
（“高校”“平台”混在一个 $N$ 里），并要求第十章
给出多群版本。本书不给完整方程（超出教学需要），
只说明多群化的两个必要改动：
其一，把 $N$ 拆为校园群与公域群，
群间接触率（弱连接强度）与群内接触率
（强连接密度）分开设参；
其二，校园群引入\keyw{学期节律项}——
寒暑假相当于大规模的暂时离场
（$S$、$I$ 的时间调制项），
它解释了北航同好会人数“新生入学后突破 200 人”
的时间点（\conf{C}）：\keyw{每年九月是一次
系统性的易感者注入}。

同时必须重申偏差警告：本章五校案例是
\keyw{被记录、被出圈}的现存样本——
全国高校中奶蛙活动存在但未被记录的院校
不可观测，扩散序列（表 \ref{tab:10:cases}）
因此是截尾样本上的形态归纳，
不是全总体估计。这与第四章“不可复核性”
（\ref{sec:04:irrevocable}）构成一对方法论提醒：
\keyw{网络文化研究的样本，一端在消失
（证据损耗），一端在缺席（未记录）}。

\section{本章小结}\label{sec:10:summary}

本章整理了五校案例（全部 \conf{C} 级），
给出扩散序列假设（采用者类别的形态相容，
\conf{M} 级）、空间实践分析（立牌作为物化符号、
流动奶蛙作为微阈限事件、FoMO 驱动纳新）、
组织形态分析（同好会的名义—实际双轨功能、
校本化表情包作为身份标记）、协商分析
（9 小时立牌事件、协商强度谱、赞诗的语体三功能）
与网络框架（强连接复读与弱连接供给的乘积、
学期节律作为易感者注入）。
全书最重的证据警告也落在本章：
所有案例均为单一来源的出圈样本，分析结论
无一例外是 \conf{M} 级诠释。遗留问题：
本章描述了协商的发生，但未测量协商的心理后果
——参与者、旁观者与管理者各自如何理解这场
“9 小时”，需要第十一章的访谈设计进入现场。

\begin{nlexercises}
  \item 用表 \ref{tab:10:spectrum} 的三维谱，
        为你所在学校设想一次假想的立牌投放事件，
        分别写出“低—中强度收尾”与“高强度升级”两个分支，
        并指出哪个环节的哪个动作是升级的转折点。
  \item “流动奶蛙合影率高于固定装置合影率”
        （\ref{sec:10:space} 的可检验推论）：
        设计一个不干预现场的观察方案
        （观察点选择、观察时段、计数规则、
        遇到的混杂因素及处理）。
  \item 本章称赞诗“用可撤回的语体，说不可撤回的态度”。
        请找一个你所见的其他社群文本（如弹幕、
        群公告、评论区接龙）检验这一概括：
        若成立，指出其可撤回装置；若不成立，
        说明该文本用了什么别的机制来处理表态风险。
\end{nlexercises}

\nlkeywords 立牌、流动奶蛙、奶 OUC、同好会、
创新扩散、采用者类别、阈限、错失恐惧（FoMO）、
弱连接、强连接、协商强度谱、低强度协商、
道德恐慌、校本化、现存样本偏差
